\chapter{Conclusions}
% What has been done with respect to what has been promised in Chapters~\ref{ch:intro} and \ref{ch:analysis}, and what is left out.

This project addressed the problem of maintaining a minimum spanning tree in a distributed
network whose topology can change over time. The implementation combines a distributed GHS
construction phase with dynamic protocols for link failures, node failures, and link
additions.

The initial MST is not supplied externally. At startup, the nodes begin as singleton
fragments and the GHS protocol constructs the initial minimum spanning forest. Once that
construction has converged, the dynamic protocols maintain the resulting tree. A tree-link
failure splits a fragment and triggers a minimum outgoing edge search, while an edge
addition identifies the heaviest edge in the resulting cycle and updates the tree when
necessary.

The project reached Tier 3 of the development plan. In particular, the implementation
includes static GHS construction, the naive and binary-search MOE procedures, addition
handling, cross-fragment reconnection, fragment re-identification, and mechanisms for
concurrent failures and additions. The pure protocol state machine can be executed both
through the OTP actor runtime and through a deterministic simulator. The validation suite
combines unit tests and simulations validated against a Kruskal-based oracle.

Tier 4 was not completed. The optional improvements proposed for that tier, such as
further optimizations and alternative distributed median-selection techniques, were left % non è molto chiaro cosa intendiamo con optimizations forse (ne qua ne nel capitolo 3)
out of the current version.

Overall, the project provides a working implementation of a dynamic distributed MST
protocol and a testing infrastructure for checking both its distributed execution and its
final global result.
